\subsection{Degradation Robustness}
\label{sec:noise}

\label{subsec:vitw-results}

Table~\ref{tab:vitw_bench} reports the sixteen-subset breakdown.
We compare AmphionASR against eight publicly reported systems
organised by access level:

\begin{itemize}
  \item \textbf{Closed-source}: Gemini~3-Flash, Seed-ASR
        \cite{bai2024seedasr}, and GPT-4o-transcribe
        \cite{openai2024gpt4o}.
  \item \textbf{Open-source}: Whisper-large-v3
        \cite{radford2023whisper}, Qwen2.5-Omni
        \cite{xu2025qwen25omni}, Kimi-Audio
        \cite{kimiaudio2025}, Qwen3-ASR-1.7B
        \cite{qwen2025qwen3asr}, and Mega-ASR
        \cite{xie2026mega} with the paper's environment-aware
        routing variant.
\end{itemize}

\begin{table}[t]
  \caption{Voices-in-the-Wild-Bench~\cite{xie2026mega} results, WER
    (\%, lower is better). The benchmark contains $5{,}000$ clips
    covering seven atomic acoustic effects plus a Mixed subset, in
    two splits: \emph{Real} (public-internet recordings, $150$ clips
    per atomic effect plus a $450$-clip Mixed) and \emph{Sim}
    (spectrogram-level simulation, $350$ clips per atomic effect plus
    a $1{,}050$-clip Mixed). Baseline numbers in the eight
    non-AmphionASR columns are reproduced from~\cite{xie2026mega}
    Table~4; the AmphionASR column is produced by the released
    checkpoint under the same normaliser as
    Section~\ref{subsec:setup}. Bold marks the lowest WER in each row.
    Baseline columns use the Mega-ASR normalisation pipeline~\cite{xie2026mega}
    rather than ours; see Section~\ref{subsec:setup} for our normaliser.}
  \label{tab:vitw_bench}
  \centering
  \scriptsize
  \setlength{\tabcolsep}{3pt}
  \begin{tabular}{l@{\hspace{3pt}}rccccccccc}
    \toprule
    & & \multicolumn{3}{c}{\emph{Closed source}} & \multicolumn{5}{c}{\emph{Open source}} & \\
    \cmidrule(lr){3-5} \cmidrule(lr){6-10}
    Subset & $n$ & Gem3-F & Seed & GPT-4o & Wh-Lv3 & Q2.5-O & Kimi-A & Q3-ASR & Mega & \textbf{Amph.} \\
    \midrule
    \multicolumn{11}{l}{\emph{Real (real-world recordings)}} \\
    \midrule
    Noise                  & 150  & 7.63  & 8.21  & 13.19 & 16.57 & 11.92 & 35.10 & 7.51  & \textbf{6.12} & 7.89 \\
    Far-field              & 150  & 5.14  & 3.06  & \textbf{1.87} & 3.38  & 2.35  & 2.71  & 2.23  & 2.33  & 2.21 \\
    Obstructed             & 150  & 3.73  & 3.10  & 1.57  & 3.06  & 2.40  & 2.49  & 1.73  & 1.80  & \textbf{1.29} \\
    Echo \& reverberation  & 150  & 8.75  & 16.55 & 15.62 & 25.34 & 20.01 & 24.00 & 10.40 & \textbf{8.66} & 8.75 \\
    Recording coloration   & 150  & 8.38  & 18.48 & 13.37 & 18.33 & 13.71 & 8.73  & 9.57  & \textbf{6.91} & 8.19 \\
    Electronic distortion  & 150  & 3.15  & 3.89  & 3.70  & 3.74  & 2.46  & 1.83  & 1.54  & 1.60  & \textbf{1.49} \\
    Transmission dropout   & 150  & 5.47  & 7.97  & 8.76  & 7.04  & 6.34  & 4.54  & 4.16  & 2.72  & \textbf{2.38} \\
    Mixed                  & 450  & 7.99  & 6.88  & 5.62  & 8.91  & 6.40  & 4.44  & 3.30  & \textbf{2.63} & 3.02 \\
    \midrule
    \multicolumn{11}{l}{\emph{Sim (spectrogram-level simulation)}} \\
    \midrule
    Noise                  & 350   & 10.61 & 8.11  & 45.78 & 18.19 & 17.88 & 14.59 & 9.52  & 8.09  & \textbf{7.79} \\
    Far-field              & 350   & 1.90  & 3.19  & 2.39  & 6.85  & 2.44  & 1.92  & \textbf{1.54} & 1.69  & 1.66 \\
    Obstructed             & 350   & 2.65  & 2.76  & 2.77  & 6.01  & 2.08  & 1.64  & \textbf{1.27} & 1.41  & 1.88 \\
    Echo \& reverberation  & 350   & 14.86 & 18.21 & 28.76 & 39.87 & 32.64 & 26.58 & 14.61 & \textbf{12.22} & 13.22 \\
    Recording coloration   & 350   & 19.85 & 23.33 & 22.60 & 31.81 & 30.09 & 18.09 & 19.42 & \textbf{13.23} & 15.08 \\
    Electronic distortion  & 350   & 7.56  & 5.71  & 8.43  & 8.77  & 5.96  & \textbf{2.78} & 3.41  & 3.35  & 3.28 \\
    Transmission dropout   & 350   & 7.65  & 7.46  & 7.71  & 8.05  & 5.88  & 6.33  & 4.19  & 2.88  & \textbf{2.50} \\
    Mixed                  & 1{,}050 & 9.62  & 9.29  & 11.00 & 14.79 & 10.29 & 6.19  & 5.39  & \textbf{4.53} & 4.87 \\
    \bottomrule
  \end{tabular}
\end{table}

\paragraph{AmphionASR leads five of sixteen subsets.}
On the \emph{Real} split, AmphionASR is the best system on the
Obstructed ($1.29$\,\% WER), Electronic distortion ($1.49$\,\%) and
Transmission dropout ($2.38$\,\%) subsets. On the \emph{Sim} split,
it leads on Noise ($7.79$\,\%) and Transmission dropout
($2.50$\,\%). The remaining eleven cells are led by Mega-ASR
(eight cells) or Qwen3-ASR-1.7B (three cells). The two systems
AmphionASR competes most closely with are therefore Mega-ASR
(the benchmark's own model, with environment-aware routing) and
Qwen3-ASR-1.7B (the closest baseline by parameter scale at
$1.7$\,B vs.\ our $4.3$\,B distinct parameters); the three
closed-source systems sit further back on most rows.

\paragraph{Compound mixtures remain the hardest condition.}
The Mixed subsets compose two to five atomic effects per clip and
are uniformly the most demanding rows of
Table~\ref{tab:vitw_bench}. AmphionASR reaches $3.02$\,\% WER on
Real-Mixed (third behind Mega-ASR's $2.63$\,\% and Qwen3-ASR-1.7B's
$3.30$\,\%) and $4.87$\,\% WER on Sim-Mixed (third behind
Mega-ASR's $4.53$\,\% and Qwen3-ASR-1.7B's $5.39$\,\%). The Real
$\!\to\!$ Sim gap on the Mixed subset is approximately $1.85$ WER
points and is broadly consistent across the nine systems,
suggesting that the spectrogram-level simulation is uniformly
slightly harder than the matched real recordings rather than
producing a system-specific gap.

\paragraph{Residual gap on echo and recording coloration.}
The two cells where AmphionASR most clearly trails Mega-ASR are
Sim-Echo \& reverberation ($13.22$ vs.\ $12.22$\,\% WER) and
Sim-Recording coloration ($15.08$ vs.\ $13.23$\,\%); the
corresponding Real-split cells are tighter ($8.75$ vs.\ $8.66$ on
Echo, $8.19$ vs.\ $6.91$ on Recording). The underlying cause has
not been investigated and a controlled decomposition by
atomic-effect type is left to the next iteration.

\paragraph{Normalisation caveat.}
Because the eight comparison columns in Table~\ref{tab:vitw_bench}
are quoted verbatim from \cite{xie2026mega} Table~4 (with the
Mega-ASR paper's normaliser) while the AmphionASR column is
produced under our own normaliser
(Section~\ref{subsec:setup}), absolute differences between the
AmphionASR column and the other eight should be read with a small
margin of system-incomparability. A like-with-like re-run of all
eight baselines through our normaliser is on the agenda for the
next revision.
